% GENERATED FILE. Edit canonical lessons, then run npm run generate:books.
% canonical-source-hash: fnv1a64:5830b003de9f0734
% canonical-lessons: PA-C114-saban, PA-C114-taulia, PA-C114-gadda, PA-C114-cadar, PA-C114-kambal

\chapter{Soap, Towel, Mattress, Bedsheet, Blanket}
\label{ch:pa-soap-towel-mattress-bedsheet-blanket}
\hlchaptermodality{114}

\begin{chapteropening}
\textbf{By the end of this chapter:} \emph{I can say soap, a towel, a mattress, a bedsheet and a blanket.}

Complete the last lesson of chapter 114: I can say soap, a towel, a mattress, a bedsheet and a blanket.
\end{chapteropening}

\section[\texorpdfstring{\pa{ਸਾਬਣ} (sābaṇ)}{ਸਾਬਣ (sābaṇ)}]{\pa{ਸਾਬਣ} (sābaṇ) — soap}

\label{lesson:PA-C114-saban}



\begin{quote}
Before the new one: say the Punjabi for a cooking pot, then the Punjabi for a comb.
\end{quote}

\subsection*{You'll want to know: \pa{ਸਾਬਣ}}
\textbf{\pa{ਸਾਬਣ}} — \emph{sābaṇ} — \textquotedblleft{}soap\textquotedblright{}.

\begin{cousinweb}
\textbf{\pa{ਸਾਬਣ}} (\emph{sābaṇ}) is a loan from Arabic \emph{ṣābūn}, \textquotedblleft{}soap\textquotedblright{}, which reached Punjabi through Persian. Its Arabic root is \emph{ṣ-b-n} (Lane).
\end{cousinweb}

\begin{grammarlens}[title={What you've built}]
One more word for this chapter.
\end{grammarlens}

\subsection*{Your turn}
\begin{itemize}
\raggedright
  \item \emph{Say it:} \emph{sābaṇ}
  \item \emph{Say it:} \emph{sābaṇ}, once more
  \item \emph{Say it:} say \emph{sābaṇ}
  \item \emph{Recall:} say the Punjabi for a cooking pot, then the Punjabi for a comb, then say \emph{sābaṇ} again
\end{itemize}

\begin{tcolorbox}[breakable,colback=teal!4,colframe=teal!35!black,title={Before you move on}]
What does \pa{ਸਾਬਣ} mean? (\textquotedblleft{}Soap\textquotedblright{}.) Say it once more.
\end{tcolorbox}
\section[\texorpdfstring{\pa{ਤੌਲੀਆ} (tauliā)}{ਤੌਲੀਆ (tauliā)}]{\pa{ਤੌਲੀਆ} (tauliā) — a towel}

\label{lesson:PA-C114-taulia}



\begin{quote}
Before the new one: say the Punjabi for a comb, then the Punjabi for soap.
\end{quote}

\subsection*{You'll want to know: \pa{ਤੌਲੀਆ}}
\textbf{\pa{ਤੌਲੀਆ}} — \emph{tauliā} — \textquotedblleft{}a towel\textquotedblright{}.

\begin{cousinweb}
\textbf{\pa{ਤੌਲੀਆ}} (\emph{tauliā}) is a loan, usually traced to Portuguese \emph{toalha}, \textquotedblleft{}towel\textquotedblright{}.
\end{cousinweb}

\begin{grammarlens}[title={What you've built}]
One more word for this chapter.
\end{grammarlens}

\subsection*{Your turn}
\begin{itemize}
\raggedright
  \item \emph{Say it:} \emph{tauliā}
  \item \emph{Say it:} \emph{tauliā}, once more
  \item \emph{Say it:} say \emph{tauliā}
  \item \emph{Recall:} say the Punjabi for a comb, then the Punjabi for soap, then say \emph{tauliā} again
\end{itemize}

\begin{tcolorbox}[breakable,colback=teal!4,colframe=teal!35!black,title={Before you move on}]
What does \pa{ਤੌਲੀਆ} mean? (\textquotedblleft{}A towel\textquotedblright{}.) Say it once more.
\end{tcolorbox}
\section[\texorpdfstring{\pa{ਗੱਦਾ} (gaddā)}{ਗੱਦਾ (gaddā)}]{\pa{ਗੱਦਾ} (gaddā) — a mattress}

\label{lesson:PA-C114-gadda}



\begin{quote}
Before the new one: say the Punjabi for soap, then the Punjabi for a towel.
\end{quote}

\subsection*{You'll want to know: \pa{ਗੱਦਾ}}
\textbf{\pa{ਗੱਦਾ}} — \emph{gaddā} — \textquotedblleft{}a mattress\textquotedblright{}.

\begin{cousinweb}
\textbf{\pa{ਗੱਦਾ}} (\emph{gaddā}) is inherited. CDIAL lists Punjabi \emph{gaddī}, \textquotedblleft{}cushion, pad, seat\textquotedblright{}, under a reconstructed \emph{*garda}, \textquotedblleft{}seat\textquotedblright{}, which it takes, probably, from Dravidian, and says Hindi \emph{gaddā}, \textquotedblleft{}quilted mattress\textquotedblright{}, was taken from Punjabi (CDIAL 4053).
\end{cousinweb}

\begin{grammarlens}[title={What you've built}]
One more word for this chapter.
\end{grammarlens}

\subsection*{Your turn}
\begin{itemize}
\raggedright
  \item \emph{Say it:} \emph{gaddā}
  \item \emph{Say it:} \emph{gaddā}, once more
  \item \emph{Say it:} say \emph{gaddā}
  \item \emph{Recall:} say the Punjabi for soap, then the Punjabi for a towel, then say \emph{gaddā} again
\end{itemize}

\begin{tcolorbox}[breakable,colback=teal!4,colframe=teal!35!black,title={Before you move on}]
What does \pa{ਗੱਦਾ} mean? (\textquotedblleft{}A mattress\textquotedblright{}.) Say it once more.
\end{tcolorbox}
\section[\texorpdfstring{\pa{ਚਾਦਰ} (cādar)}{ਚਾਦਰ (cādar)}]{\pa{ਚਾਦਰ} (cādar) — a bedsheet}

\label{lesson:PA-C114-cadar}



\begin{quote}
Before the new one: say the Punjabi for a towel, then the Punjabi for a mattress.
\end{quote}

\subsection*{You'll want to know: \pa{ਚਾਦਰ}}
\textbf{\pa{ਚਾਦਰ}} — \emph{cādar} — \textquotedblleft{}a bedsheet\textquotedblright{}.

\begin{cousinweb}
\textbf{\pa{ਚਾਦਰ}} (\emph{cādar}) is a loan, usually traced to Persian \emph{chādar}, \textquotedblleft{}sheet, veil\textquotedblright{}.
\end{cousinweb}

\begin{grammarlens}[title={What you've built}]
One more word for this chapter.
\end{grammarlens}

\subsection*{Your turn}
\begin{itemize}
\raggedright
  \item \emph{Say it:} \emph{cādar}
  \item \emph{Say it:} \emph{cādar}, once more
  \item \emph{Say it:} say \emph{cādar}
  \item \emph{Recall:} say the Punjabi for a towel, then the Punjabi for a mattress, then say \emph{cādar} again
\end{itemize}

\begin{tcolorbox}[breakable,colback=teal!4,colframe=teal!35!black,title={Before you move on}]
What does \pa{ਚਾਦਰ} mean? (\textquotedblleft{}A bedsheet\textquotedblright{}.) Say it once more.
\end{tcolorbox}
\section[\texorpdfstring{\pa{ਕੰਬਲ} (kambal)}{ਕੰਬਲ (kambal)}]{\pa{ਕੰਬਲ} (kambal) — a blanket}

\label{lesson:PA-C114-kambal}



\begin{quote}
Before the new one: say the Punjabi for a mattress, then the Punjabi for a bedsheet.
\end{quote}

\subsection*{You'll want to know: \pa{ਕੰਬਲ}}
\textbf{\pa{ਕੰਬਲ}} — \emph{kambal} — \textquotedblleft{}a blanket\textquotedblright{}.

\begin{cousinweb}
\textbf{\pa{ਕੰਬਲ}} (\emph{kambal}) is a tadbhava, a word inherited from Sanskrit and reshaped by sound change. It comes from \emph{kambala}, \textquotedblleft{}woollen blanket\textquotedblright{} (CDIAL 2771). CDIAL notes a possible Austroasiatic origin, and says Hindi \emph{kambal} was taken from Punjabi.
\end{cousinweb}

\begin{grammarlens}[title={What you've built}]
One more word for this chapter.
\end{grammarlens}

\subsection*{Your turn}
\begin{itemize}
\raggedright
  \item \emph{Say it:} \emph{kambal}
  \item \emph{Say it:} \emph{kambal}, once more
  \item \emph{Say it:} say \emph{kambal}
  \item \emph{Recall:} say the Punjabi for a mattress, then the Punjabi for a bedsheet, then say \emph{kambal} again
\end{itemize}

\begin{tcolorbox}[breakable,colback=teal!4,colframe=teal!35!black,title={Before you move on}]
What does \pa{ਕੰਬਲ} mean? (\textquotedblleft{}A blanket\textquotedblright{}.) Say it once more.
\end{tcolorbox}
